como los numeros complejos con finalmente numeros en R x R 
plano de argande:

coordenas polares 

trazo la recta que va del origen al punto (x,y) y esa distancia es r
tambien puedo calcularlo el angulo theta 

tg $\theta$ = y/x 
sen $\theta$ = y/r
cos $\theta$ = x/r

como se que el arcotangente de theta me da el theta? 
depende de las coordenadas de (x,y)

676767676767676767676767
a>1 r,s in Q
n <= m -> a^n a^m 

z = r e^itheta = r cos thetha + i sen thetha
w = s e^ialfa = s(cos alfa + i sen alfa)

teorema de euler
e^itetha =  cos theta + i sen theta
zw = r(cos tetha + i sen theta) * s(cos alfa + i sen alfa)
rs(cos o con alf - sen 0 sen alf) + i (cos 0 sen alf + sen 0 cos alf)
rs ( cos(0+alf) + isen(0+alf))

rs * e^i(0+alf)
arg(zw) = arg(z)+arg(w)
z tiene varios args??

++++++++++++++++++++++++++++++++++++++++++++++++++++++++++++++++++++++++++++++++++++++++++++++++++

(cos 0 + isen)elevado n =  z ^ n = e^in0 = cos 0n + isen 0n

z = r e ^i0  , w^2 = z 
sup que w = S e^ietha . lo que quioer ver que S^2 * e^i2etha = w^2 = z = re^io
2etha - tetha = 2kpi

S = raiz cuadrada r 
etha = theta/ 2 + pik  / k in Z

k= 0 etha = thetha/2
k= 1 etha = theta / 2 +pi
k=2  etha = theta / 2 +2pi
se vuelve a repetir  

los numeros impares n.esja de orden par
todo numero complejo tiene exactamente n-raices e-nesimas , que son 

$\raiz enesima de r * e^io/n , raiz enemsima de r * e^i/thata/n *2pi/n$


raiz enesima de r * e^i(0/n * k 2pi/n) = CUIDADO CALAVERITA = raiz enesima Z 

1 = 1 * e^iO 
las raices nesimas de 1 son los numeros de la forma :
e^i2pik/n 
def : 
Gn = { w in C : w^n = 1 ^ n}
tema : |G | = n 
Gn =  {e^i2pik/n} 

es cerrado para el producto , si uno tiene dos raices para la unidad
si z , w in Gn entonces zw in Gn 
z^n = 1 = w ^ n -> (zw)^n = z^n * w^n = 1*1 = 1 
(z^-1)^n = z^-n = 1/1 = 1 
1 in Gn 
Gn es un grupo Abeliano para el producto 
entonces , Gn <= Gm , z in Gn z^n = 1 z^2n = (z^n)^2 = 1^2 = 1 , si y solo si n | m

si Gn <= Gm entonces n|m , 
si que 1 = (e^i2pi/n)^m = e^i2pim/n
entonces hay un entero tal que 2pi - 0 = 2kpi
m/n = k  -> n = k*m asi que n|m